\section{Related Work}
\label{sec:related-work}

\parbench{} draws on three lines of work: code translation benchmarks, LLM-for-HPC evaluation, and robustness testing.

\paragraph{General-purpose code benchmarks.}
Execution-based evaluation is now the standard for measuring functional correctness, but existing task distributions do not cover parallel HPC translation. HumanEval~\cite{HumanEval2021} and TransCoder~\cite{TransCoder2020} operate at function granularity, while SWE-bench~\cite{SWEbench2024} evaluates repository-level issue resolution in general software projects. These benchmarks do not exercise GPU memory models, thread synchronization, API-specific launch structure, or host-device coordination. SWE-bench Illusion~\cite{SWEbenchIllusion2025} further demonstrates that static benchmark scores can overstate model capability when evaluation methodology is not carefully controlled, a concern that motivates \parbench{}'s conjunctive build-run-verify pipeline and augmentation robustness probes.

\paragraph{Parallel-code benchmarks.}
ParEval~\cite{ParEval2024} studies synthesis of parallel code from natural language tasks, testing whether a model can produce a solution but not whether it can preserve an existing implementation's decomposition across APIs. ParEval-Repo~\cite{ParEvalRepo2025} studies full-repository translation and demonstrates a severe integration bottleneck: models fail once they must recreate build systems and repository scaffolding. \parbench{} is complementary to both, it preserves executable evaluation but fixes build infrastructure so the measured object is kernel translation. VibeCodeHPC~\cite{VibeCodeHPC2025} evaluates prompt-driven HPC code generation, while QiMeng-MuPa~\cite{QiMengMuPa2025} addresses multi-paradigm parallel code synthesis; both focus on generation rather than translation of existing parallel implementations.

\paragraph{HPC-specific LLM evaluation.}
LASSI~\cite{LASSI2024} is the nearest prior work in granularity and verification. It evaluates an agentic self-correction pipeline on 10 HeCBench kernels across two directions (CUDA to OpenMP-target and reverse), measuring how iterative repair improves translation success. \parbench{} asks a different question - what is the first-attempt translation capability across a broad set of directions? 
and differs in scope: 35~kernels from five suites, six standard bidirectional directions among CUDA, OpenMP, and OpenCL plus four OpenMP-target case-study directions, AST-driven augmentation for robustness testing, and a declarative spec format designed for community extension beyond the initial corpus. CodeRosetta~\cite{CodeRosetta2024} and HPC-Coder-V2~\cite{HPCCoderV2} demonstrate specialized modeling and fine-tuning for parallel code, but \parbench{} provides an execution-based substrate for comparing such models against general-purpose LLMs. Extended comparisons are in Appendix~\ref{sec:appendix-related}.

\paragraph{Benchmark validity and robustness.}
Public code benchmarks can be memorized, and surface-form perturbations can reveal when evaluation depends on brittle pattern matching~\cite{SWEbenchIllusion2025}. \parbench{}'s augmentation engine is not a complete solution to training-data contamination because it cannot rule out algorithm-level familiarity, but it provides a concrete robustness probe for widely circulated HPC kernels. Its correctness-only scope also differs from efficiency-oriented work such as TRACY~\cite{TRACY2025}, which shows that functional correctness and performance are separable.
